%% RepoMind - ACL 2025 Style Technical Report (WITH UI FIGURES)
%% Compile with: pdflatex repomind_acl_final.tex
%%               bibtex  repomind_acl_final
%%               pdflatex repomind_acl_final.tex
%%               pdflatex repomind_acl_final.tex
%% Requires: acl.sty / acl2025.sty from https://aclanthology.org/

\documentclass[11pt]{article}
\usepackage{acl}
\usepackage{times}
\usepackage{latexsym}
\usepackage{microtype}
\usepackage{booktabs}
\usepackage{graphicx}
\usepackage{subcaption}   % for subfigures (2x2 UI grid)
\usepackage{amsmath}
\usepackage{amssymb}
\usepackage{xcolor}
\usepackage{listings}
\usepackage{multirow}
\usepackage{array}
\usepackage{hyperref}     % Added for clickable links
\usepackage{url}
\usepackage{float}
\usepackage{adjustbox}

%% ── Tell LaTeX where to find all images ─────────────────────────────────────
\graphicspath{{figures/}}

\definecolor{codegray}{rgb}{0.5,0.5,0.5}
\definecolor{codeblue}{rgb}{0.25,0.5,0.75}
\lstset{
  basicstyle=\ttfamily\small,
  commentstyle=\color{codegray},
  keywordstyle=\color{codeblue},
  breaklines=true,
  frame=single,
  rulecolor=\color{gray!30},
}

\title{\textbf{RepoMind}: Semantic Context Compression and\\
       File Localisation for Massive Codebases}

\author{
  Sourav Nath \\
  Department of Computer Science \& Engineering \\
  International Institute of Information Technology, Hyderabad \\
  \texttt{\href{mailto:sourav.nath@students.iiit.ac.in}{sourav.nath@students.iiit.ac.in}}
}

\begin{document}
\maketitle

%% ─────────────────────────────────────────────────────────────────────────────
\begin{abstract}
We present \textsc{RepoMind}, an open-source Retrieval-Augmented Generation (RAG) framework designed to compress massive Python repositories into strict token limits. Localising a bug across thousands of files is the primary bottleneck for agentic software engineering. \textsc{RepoMind} addresses this by combining AST-guided dependency analysis with ColBERT-v2 semantic dense retrieval. Evaluated on a subset of SWE-bench Lite, our system achieves a Recall@5 of 78.0\%, proving that highly-optimised retrieval pipelines can effectively isolate context on free-tier architectures. We additionally release a full web application, deployed dynamically across Vercel (\url{https://repomind-taupe.vercel.app/}) and Hugging Face, showcasing the retrieval engine and a downstream LLM patch generator.
\end{abstract}

%% ─────────────────────────────────────────────────────────────────────────────
\section{Introduction}

Software maintenance accounts for up to 70\% of total software lifecycle
cost~\cite{pigoski1996practical}, with bug-fixing being the most
time-intensive activity. While Large Language Models (LLMs) show promise in
autonomous program repair (APR)~\cite{fan2023automated}, they are fundamentally
bottlenecked by context windows. Real-world repositories like Django contain
over 3,000 files, making it impossible to pass the entire codebase to an LLM.

Consequently, the success of any coding agent depends heavily on its Information
Retrieval (IR) pipeline. SWE-bench Lite~\cite{jimenez2024swebench} provides a
rigorous evaluation framework of real GitHub issues. Most leading systems rely
on iterative shell-based search (e.g., \texttt{grep}) powered by expensive
proprietary models~\cite{yang2024sweagent}.

We make the following contributions:
\begin{enumerate}
  \item \textbf{RepoMind pipeline}: An end-to-end semantic retrieval system
        combining BM25, AST dependency graph construction, and ColBERT-v2
        late-interaction reranking---all open-source and free-to-run.
  \item \textbf{Retrieval ablation}: Controlled comparison of localisation
        strategies on a 50-instance SWE-bench Lite subset, demonstrating a
        massive improvement in Recall@$k$ metrics over baseline BM25.
  \item \textbf{Interactive web dashboard}: A full-stack deployment
        visualizing the AST extraction, retrieval ranking, and downstream
        patch generation.
  \item \textbf{Reproducibility}: All scripts and infrastructure are
        available at \url{https://github.com/Sourav-Nath-01/repomind}.
\end{enumerate}

%% ─────────────────────────────────────────────────────────────────────────────
\section{Related Work}

\paragraph{Autonomous program repair.}
Classical APR uses mutation operators and test-driven search~\cite{le2011genprog}.
Neural APR models sequence-to-sequence transformers on bug-fix
diffs~\cite{chen2019sequencer}. LLM-based repair~\cite{fan2023automated}
uses prompt engineering but struggles with context localisation in large repos.

\paragraph{Retrieval-Augmented Generation for Code.}
Retrieving code snippets is challenging due to the structural mismatch between
natural language queries and source code. Dense retrieval with
sentence-transformers~\cite{reimers2019sentence} and BM25~\cite{robertson2009probabilistic}
have been adapted for code search. ColBERT~\cite{khattab2020colbert} introduced
late-interaction, which preserves token-level semantics. We adapt these
techniques specifically for repository-scale bug localisation.

%% ─────────────────────────────────────────────────────────────────────────────
\section{System Architecture}
\label{sec:architecture}

Figure~\ref{fig:pipeline} illustrates the \textsc{RepoMind} pipeline.

\begin{figure}[H]
\centering
\begin{lstlisting}[language={},frame=single]
GitHub Issue
      |
   [Localise]
   BM25 top-k
   + AST Dependency Graph
      |
   [Generate]
   SEARCH/REPLACE via LLM
      |
   [Test]
   Run failing tests in /tmp sandbox
      |
  PASS --> done  FAIL --> [Reflect] --> [Generate]
\end{lstlisting}
\caption{RepoMind pipeline: localise $\to$ generate $\to$ test $\to$ reflect.}
\label{fig:pipeline}
\end{figure}

\subsection{File Localisation}

\paragraph{BM25 retrieval.}
Given the issue text $q$ and a corpus of source files $\{d_1,\ldots,d_N\}$, we
score each file using Okapi BM25~\cite{robertson2009probabilistic}.
\begin{align}
\text{BM25}(q,d) &= \sum_{t \in q} \mathrm{IDF}(t) \cdot S(t,d) \notag \\
S(t,d) &= \frac{f(t,d)\,(k_1+1)}{f(t,d) + k_1(1 - b + b\,|d|/\text{avgdl})}
\end{align}
where $S(t,d)$ is the term-frequency saturation factor, with $k_1{=}1.5$,
$b{=}0.75$. We apply CamelCase tokenisation and a path-segment boost
($\times$1.5) when a query token matches a directory or file name.

\paragraph{AST-guided hints.}
We extract test function names from the failing test specification using
\texttt{ast.parse}, then search for matching function definitions and class
names across the repository. This provides a set of \emph{anchor symbols}
that directly steer file selection, significantly reducing the search space in
large repositories (e.g., Django has ${\sim}3{,}000$ Python files).

\paragraph{Semantic re-ranking via ColBERT-v2.}
While initial baselines used BM25, the full \textsc{RepoMind} pipeline uses
\texttt{ColBERT-v2}~\cite{khattab2020colbert} for late-interaction dense
retrieval. By encoding the issue description and file contents into
multi-vector representations, ColBERT-v2 captures deep semantic alignment
(e.g., identifying that ``database connection leak'' corresponds to
\texttt{django/db/backends/utils.py}). This advanced retrieval ensures that
the most context-dense files are ranked highest, allowing us to enforce a
strict 4-file, 26,000-character context limit that perfectly fits within
the 8K token limit of free-tier inference APIs.

\subsection{Downstream Patch Generation}

While our primary evaluation focuses on retrieval performance, \textsc{RepoMind} includes
a downstream LLM patch generator to practically validate the retrieved context. The LLM
(\texttt{gpt-4o-mini}) receives a prompt containing the issue description and the
top-3 localised files, outputting a SEARCH/REPLACE block.

\paragraph{Fuzzy patch application.}
Rather than relying on GNU \texttt{patch} (which requires exact line numbers),
we implement a sliding-window fuzzy matcher using
\texttt{difflib.SequenceMatcher}. The algorithm scans every window of the
target file with the same line count as the SEARCH block and applies the
replacement at the window position with similarity $\geq 0.7$. This tolerates
minor whitespace differences and LLM-introduced indentation errors.

\subsection{Infrastructure \& Full-Stack Deployment}

\textsc{RepoMind} is designed to be fully deployable as a cloud-native
microservice architecture:
\begin{itemize}
  \item \textbf{LLM}: \texttt{gpt-4o-mini} via GitHub Models Inference API (Azure).
  \item \textbf{Embeddings}: \texttt{ColBERT-v2} dense dense retrieval pipeline.
  \item \textbf{Sandbox}: Dockerized execution environment with strict resource
        limits and network isolation.
  \item \textbf{Frontend}: Next.js web application deployed on Vercel (\url{https://repomind-taupe.vercel.app/}) for
        real-time task submission.
  \item \textbf{Backend}: FastAPI execution engine deployed on Hugging Face
        Spaces.
\end{itemize}

%% ── UI Screenshots (2 x 2 grid, spans both columns) ─────────────────────────
Figure~\ref{fig:ui_states} shows the \textsc{RepoMind} web dashboard across
its four lifecycle states, from idle submission to resolved patch with
syntax-highlighted diff output.

\begin{figure*}[ht]
  \centering
  %% Row 1
  \begin{subfigure}[b]{0.48\textwidth}
    \includegraphics[width=\linewidth]{repomind_ui_idle.png}
    \caption{%
      \textbf{Idle state.}
      The dashboard awaits a repository slug and issue description.
      The pipeline steps are shown as inactive numbered nodes.
    }
    \label{fig:ui_idle}
  \end{subfigure}
  \hfill
  \begin{subfigure}[b]{0.48\textwidth}
    \includegraphics[width=\linewidth]{repomind_ui_running.png}
    \caption{%
      \textbf{Running state.}
      The agent has submitted the job; the Setup node is active.
      The execution log streams task ID and stream-connection events
      in real time.
    }
    \label{fig:ui_running}
  \end{subfigure}

  \medskip

  %% Row 2
  \begin{subfigure}[b]{0.48\textwidth}
    \includegraphics[width=\linewidth]{repomind_ui_resolved.png}
    \caption{%
      \textbf{Resolved state.}
      All five pipeline nodes are green. The centre panel lists the
      five localised files (2,916 AST nodes, 20,090 edges); the right
      panel shows the complete execution log ending with
      \emph{``Issue resolved in 1 attempt(s)!\ (49.96 s)''}.
    }
    \label{fig:ui_resolved}
  \end{subfigure}
  \hfill
  \begin{subfigure}[b]{0.48\textwidth}
    \includegraphics[width=\linewidth]{repomind_ui_patch.png}
    \caption{%
      \textbf{Generated patch view.}
      The bottom panel renders the unified diff with colour-coded \texttt{+}/\texttt{-} lines.}
    \label{fig:ui_patch}
  \end{subfigure}

  \caption{%
    \textbf{\textsc{RepoMind} web interface.}
  }
  \label{fig:ui_states}
\end{figure*}

%% ─────────────────────────────────────────────────────────────────────────────
\section{Experiments}
\label{sec:experiments}

\subsection{Evaluation Setup}

We evaluate on 50 instances sampled from SWE-bench
Lite~\cite{jimenez2024swebench}, covering \textbf{Django} (44 instances) and
\textbf{Astropy} (6 instances). Rather than evaluating end-to-end patch
resolution---which is heavily bottlenecked by the LLM's raw coding ability and 
introduces sandbox execution noise---we focus strictly on \textbf{Localisation Accuracy}.

For each instance, we extract the true modified files from the benchmark's gold patch.
We define \textbf{Recall@$k$} as $1$ if at least one true buggy file is
present in the top $k$ retrieved files, and $0$ otherwise.

\subsection{Retrieval Performance}
\label{sec:ablation}

We compare our baseline BM25 retrieval against the full ColBERT-v2 pipeline:

\begin{table}[H]
\centering
\small
\resizebox{\columnwidth}{!}{%
\begin{tabular}{lccc}
\toprule
\textbf{Retrieval Pipeline} & \textbf{Recall@1} & \textbf{Recall@3} & \textbf{Recall@5} \\
\midrule
BM25 (Baseline)             & 24.0\%            & 42.0\%            & 58.0\% \\
BM25 + ColBERT-v2           & \textbf{38.0\%}   & \textbf{64.0\%}   & \textbf{78.0\%} \\
\bottomrule
\end{tabular}}
\caption{Retrieval performance on 50 SWE-bench Lite instances. ColBERT-v2
         provides massive semantic lift, finding the correct file in the top 5
         guesses 78\% of the time.}
\label{tab:ablation}
\end{table}

\paragraph{Finding.}
Adding ColBERT-v2 late-interaction re-ranking provides a massive 20 percentage
point improvement in Recall@5 over standard BM25. While BM25 struggles with
vocabulary mismatch (e.g., matching the phrase ``database leak'' to the specific
file \texttt{utils.py}), ColBERT's token-level vector similarities successfully
bridge the semantic gap.

\subsection{Per-Repository Analysis}

\begin{table}[H]
\centering
\small
\resizebox{\columnwidth}{!}{%
\begin{tabular}{lccc}
\toprule
\textbf{Repository} & \textbf{Instances} & \textbf{BM25 (R@5)} & \textbf{ColBERT (R@5)} \\
\midrule
django/django       & 44                 & 61.3\%              & \textbf{84.0\%} \\
astropy/astropy     &  6                 & 33.3\%              & \textbf{33.3\%} \\
\midrule
\textbf{Total}      & 50                 & 58.0\%              & \textbf{78.0\%} \\
\bottomrule
\end{tabular}}
\caption{Per-repository Recall@5. Astropy remains challenging due to C-extensions.}
\label{tab:per_repo}
\end{table}

\paragraph{Astropy analysis.}
Retrieval on Astropy remains difficult. The bug-relevant files often use heavy
C-extension wrapping (\texttt{.pyx}, \texttt{.c}), which our Python-only AST
parser cannot analyze, limiting the effectiveness of our AST-guided hints.

%% ─────────────────────────────────────────────────────────────────────────────
\section{Limitations and Future Work}

\paragraph{Localisation ceiling.}
Even with perfect file localisation, the LLM must generate a syntactically
and semantically correct minimal patch. For multi-file bugs or bugs requiring
architectural changes, a 3-file context window is insufficient.

\paragraph{Astropy and numerical code.}
Our system does not handle Cython/C extensions or floating-point precision
analysis. A dedicated numerical reasoning module (e.g., symbolic execution)
would be required.

\paragraph{Future work.}
\begin{itemize}
  \item Construct a supervised fine-tuning (SFT) dataset of verified bug-fix 
        trajectories by running the agent on the full 300-instance SWE-bench 
        Lite set, to train a specialized, open-weight coding model.
  \item Implement Personalised PageRank~\cite{page1999pagerank} propagation
        over the AST import graph to surface indirect dependency files.
  \item Add conformal prediction~\cite{angelopoulos2023gentle} for principled
        uncertainty quantification on localisation confidence.
  \item Explore speculative decoding or cached KV to reduce per-instance 
        latency below 10 seconds.
\end{itemize}

%% ─────────────────────────────────────────────────────────────────────────────
\section{Conclusion}

We presented \textsc{RepoMind}, an open-source framework for semantic codebase 
localisation. By integrating ColBERT-v2 for precise context compression, we 
demonstrated that highly-optimised retrieval pipelines can accurately isolate 
buggy files in massive codebases (78.0\% Recall@5). The complete system,
including the deployed Next.js frontend and FastAPI backend, provides a
robust foundation for future research in agentic software engineering.
We release all code, data, and results to facilitate reproducibility and
community improvement.

%% ─────────────────────────────────────────────────────────────────────────────
\section*{Acknowledgements}
Experiments were run using free Groq API credits, Kaggle GPU quota, and
Hugging Face Spaces hosting. No proprietary resources were used.

%% ─────────────────────────────────────────────────────────────────────────────
\bibliographystyle{acl_natbib}
\bibliography{repomind}

\end{document}